\chapter{从单片机裸机程序开始}

\section{原理}

理解操作系统，最稳的起点不是 Linux，也不是进程和虚拟内存，而是一块最小单片机。先看一个具体任务：做一块温控板。

这块板子要做四件事：

\begin{enumerate}
\tightlist
\item 每 1ms 读取温度传感器。
\item 每 1ms 根据温度更新电机或风扇输出。
\item 随时接收串口命令，例如修改目标温度。
\item 每 100ms 通过串口发一行日志。
\end{enumerate}

第一版最容易写成这样：

\begin{lstlisting}
init_clock();
init_sensor();
init_motor();
init_uart();

while (true) {
  read_sensor();
  update_motor();
  handle_uart_command_if_any();
  send_log_if_due();
}
\end{lstlisting}

这就是裸机程序：没有操作系统，程序直接拥有整台机器。所有内存都能读写，所有外设寄存器都能访问，任何函数都可以关中断、改时钟、写串口、改 GPIO。它的优点是简单、启动快、路径短；缺点也同样直接：只要某个函数等硬件，后面的函数就全被拖住。

假设 \code{send\_log\_if\_due()} 里要发 80 个字节，而串口暂时发送缓冲区满。最朴素写法会忙等：

\begin{lstlisting}
void uart_putc(char c) {
  while (!uart_tx_ready()) {
    // wait until hardware can accept one byte
  }
  uart_write(c);
}
\end{lstlisting}

如果这一等花了 50ms，那么 1ms 周期的 \code{read\_sensor()} 和 \code{update\_motor()} 就会迟到 50 次。程序不是“逻辑错”，而是结构错：它把需要等待的 I/O 当成了立刻完成的普通函数。

所以本章的过渡线不是从术语开始，而是从这个失败开始：

\begin{enumerate}
\tightlist
\item 暴力写法：所有工作按顺序塞进一个 \code{while (true)}。
\item 发现问题：相邻工作共享同一个 CPU，慢 I/O 会拖死周期任务。
\item 第一步优化：把任务的周期、下次运行时间和状态显式记录下来。
\item 第二步优化：把阻塞发送改成状态机，硬件没准备好就先返回。
\item 再往后：让硬件用中断通知事件，把“主动一直问”改成“事件到了再处理”。
\end{enumerate}

\subsection{一个 MCU 芯片里面通常有什么}

现在再看硬件部件。单片机，也叫 MCU，是一颗芯片上集成了一套很小的计算机。温控板不需要插显卡和硬盘，但它需要 CPU 核执行指令，需要 Flash 保存程序，需要 SRAM 放栈和全局变量，需要 GPIO 控制电机引脚，需要 UART 收发命令，需要定时器提供 1ms 节拍。很多外设不是运行时枚举出来，而是在芯片手册里固定了基地址和寄存器布局。

\begin{longtable}{p{0.2\linewidth}p{0.34\linewidth}p{0.34\linewidth}}
\toprule
部件 & 作用 & 与 OS 的关系 \\
\midrule
CPU core & 执行指令和异常入口 & 上下文、特权模式、中断响应 \\
Flash & 保存程序镜像和只读数据 & 类似启动介质和只读代码段 \\
SRAM & 保存栈、堆、全局变量 & 物理内存管理的最小形态 \\
中断控制器 & 管理外设中断 & IRQ 分发和优先级的原型 \\
SysTick/Timer & 周期时间基准 & tick、超时和抢占的原型 \\
GPIO、UART、SPI、I2C & 连接外部设备 & 字符设备和总线驱动的原型 \\
DMA controller & 在外设和内存间搬运数据 & DMA 生命周期和 completion 的原型 \\
\bottomrule
\end{longtable}

这些部件通常共享同一个地址空间。CPU 访问 0x20000000 可能是 SRAM，访问 0x40000000 可能是 UART 寄存器。区别不在指令，而在地址译码后连到哪个硬件。地址译码可以理解为“硬件看地址范围，把这次访问交给内存或某个外设”。这个事实会一路延伸到大型 OS：驱动通过 MMIO 地址访问 PCIe 设备寄存器，本质仍是 load/store 触发硬件状态机。

\subsection{芯片设计者怎样选择寄存器接口}

外设寄存器是硬件和软件的协议。设计一个 UART 外设时，芯片设计者通常会暴露控制寄存器、状态寄存器、数据寄存器、波特率寄存器和中断使能寄存器。软件写控制寄存器启用发送接收；读状态寄存器判断 TX ready 或 RX available；读写数据寄存器收发字节。

\begin{lstlisting}
UART register model:
  CTRL     bit0 enable_rx, bit1 enable_tx
  STATUS   bit0 rx_ready, bit1 tx_empty, bit2 overrun
  DATA     read pops received byte, write submits transmit byte
  BAUD     divisor for baud rate generator
  IRQ_EN   enable rx/tx/error interrupt
\end{lstlisting}

注意 \code{DATA} 寄存器不是普通内存。读它可能弹出接收 FIFO 的一个字节，写它可能启动发送状态机。状态位也可能有特殊清除规则，例如“写 1 清除”。驱动必须严格按手册顺序访问这些寄存器，否则问题不是编译错误，而是硬件进入不可预期状态。

\subsection{从外设状态机看驱动}

一个 UART 发送器可以看成硬件状态机：空闲时等待软件写 DATA；收到字节后按波特率移位输出；完成后设置 tx\_empty；若开启中断，则向中断控制器发请求。软件驱动的任务不是“调用发送函数”，而是和这个状态机同步。

\begin{lstlisting}
UART_TX_IDLE:
  if software writes DATA:
    shift_reg = DATA
    state = UART_TX_BUSY

UART_TX_BUSY:
  shift one bit per baud tick
  if all bits sent:
    STATUS.tx_empty = 1
    if IRQ_EN.tx_empty:
      raise_irq()
    state = UART_TX_IDLE
\end{lstlisting}

这就是为什么驱动要处理 ready 位、超时、FIFO 满、错误位和中断确认。ready 位表示“硬件现在能不能接收下一步动作”；FIFO 是硬件里的小队列，用来暂存还没处理完的字节；错误位记录溢出、校验失败或线路异常。硬件不会执行 C 函数返回值协议，它只会按寄存器位和状态机推进。操作系统把这种硬件状态机包装成 \code{read/write/poll/ioctl}，但底层仍然必须尊重硬件协议。

\subsection{暴力写法：一个主循环包办所有工作}

\begin{lstlisting}
init_clock();
init_gpio();
init_uart();

while (true) {
  if (uart_has_byte()) {
    byte b = uart_read();
    handle_command(b);
  }
  if (timer_expired()) {
    update_control_loop();
  }
  flush_logs_if_needed();
}
\end{lstlisting}

这段代码看起来不像 OS，但它已经暴露了 OS 的胚胎问题。第一，多个活动共享一个 CPU，谁先执行？第二，事件到达时间不可控，怎样避免漏事件？第三，某个处理函数太慢，怎样不拖死其他工作？第四，共享状态由谁保护？第五，故障后怎样恢复到可解释状态？

注意这里的“任务”还不是 Linux 里的进程或线程。它只是一个需要周期性或事件性执行的函数。先把这个简单含义抓住，后面再逐步升级成调度器里的 task。

\subsection{从超级循环到任务表}

裸机主循环通常被称为超级循环。它的结构是一个大循环里按固定顺序调用多个任务函数。温控板的第一版可以直接写：

\begin{lstlisting}
while (true) {
  task_read_sensor();
  task_control_motor();
  task_send_telemetry();
  task_blink_led();
}
\end{lstlisting}

这种写法的问题是隐含了调度策略。调度策略就是“谁先用 CPU，谁后用 CPU，最多能用多久”的规则。超级循环没有把规则写成数据；它把规则藏在代码顺序里：每个任务的优先级就是它在代码里的顺序，每个任务的时间预算完全靠自觉，任何任务卡住，后面的任务都不会运行。

要让它可分析，第一步是把任务放进表里，并给每个任务明确的周期、deadline 和上次运行时间。deadline 是“最晚必须完成的时间点”。例如控制电机每 1ms 必须更新一次，超过这个时间点，系统还能跑，但控制效果已经变差。

\begin{lstlisting}
struct Task {
  const char* name;
  uint32_t period_ms;
  uint32_t next_release_ms;
  void (*run)();
};

while (true) {
  uint32_t now = ticks_ms();
  for each task in tasks:
    if (now >= task.next_release_ms) {
      task.run();
      task.next_release_ms += task.period_ms;
    }
}
\end{lstlisting}

这个算法是最小合作式调度器。合作式的意思是：任务必须主动返回，系统才有机会运行下一个任务。它没有抢占；抢占是“任务还没主动返回，定时器或内核入口就把 CPU 收回来”。这里先不要急着上抢占，先看合作式的边界。

它的复杂度是每轮扫描 \code{O(n)}，n 是任务数量；它的关键不变量是 \code{next\_release\_ms} 单调前进，每个任务函数必须短小、不能无限循环、不能执行不可控阻塞。

合作式调度器已经比手写超级循环好，因为任务的释放条件变成数据。但它仍然无法处理一个任务超时的问题。如果 \code{task\_send\_telemetry} 等串口缓冲区空等了 50ms，那么 1ms 周期的控制任务就会迟到。于是我们需要区分“能立刻做的计算”和“需要等待的 I/O”，这会把系统推向事件队列和中断。

\subsection{最小调度算法的复杂度和失败边界}

周期任务表的扫描算法很容易写，但它把每轮成本变成 \code{O(n)}。如果任务数量很少，成本可忽略；如果任务很多或周期很短，调度器本身会吃掉可观 CPU 时间。更重要的是，它没有强制时间预算，只能在任务返回后继续扫描。

\begin{longtable}{p{0.2\linewidth}p{0.34\linewidth}p{0.34\linewidth}}
\toprule
策略 & 算法成本 & 失败边界 \\
\midrule
固定顺序超级循环 & 每轮 \code{O(n)} & 任务顺序即隐式优先级，慢任务拖累后续任务 \\
周期任务表 & 每轮扫描 \code{O(n)} & 任务不返回则全局停顿，deadline 只能事后发现 \\
优先级就绪队列 & 取最高优先级可为 \code{O(1)} 或 \code{O(log n)} & 低优先级可能饥饿，需要抢占和时间统计 \\
时间片轮转 & 每次切换接近 \code{O(1)} & 交互延迟和吞吐受时间片影响 \\
\bottomrule
\end{longtable}

这里已经能看见调度器设计的核心：策略不是口号，必须落到数据结构和复杂度。任务少时，简单扫描最可靠；任务多、优先级复杂、需要抢占时，就要引入 runqueue、时间统计和上下文切换。

\subsection{状态机替代阻塞函数}

裸机程序最常见的灾难是阻塞等待：

\begin{lstlisting}
void send_packet(Packet p) {
  for each byte in p:
    while (!uart_tx_ready()) {
      // busy wait
    }
    uart_write(byte);
}
\end{lstlisting}

这段代码在功能上正确，但会占住整个 CPU。它的问题和上一章滑动窗口里“每次重新统计整个窗口”类似：明明只有“串口暂时不能收下一个字节”这一点新情况，程序却把整个 CPU 都扣在原地。更好的写法是把发送过程改成状态机：如果硬件暂时不能发送，就保存进度并返回，让其他任务继续运行。

\begin{lstlisting}
enum TxState { Idle, Sending };

void uart_tx_step() {
  if (tx_state == Idle || !uart_tx_ready()) {
    return;
  }
  uart_write(tx_buffer[tx_pos++]);
  if (tx_pos == tx_len) {
    tx_state = Idle;
  }
}
\end{lstlisting}

状态机的核心代价是复杂性。阻塞函数把“下一步在哪里”藏在调用栈里；状态机把下一步显式放在 \code{tx\_state} 和 \code{tx\_pos} 里。操作系统中的线程，就是另一种保存“下一步在哪里”的方法：线程把寄存器、栈和程序计数器保存为上下文，使阻塞代码看起来仍像顺序代码。但线程要付出调度和栈内存成本。

\subsection{资源所有权和最小内核对象}

没有 OS 时，所有模块都能直接操作硬件寄存器。小程序可以这样写；程序变大后，它会造成所有权混乱。串口驱动、日志模块、命令行模块如果都直接写 UART 寄存器，就无法保证输出顺序，也无法处理并发写入。

于是我们引入最小内核对象：一个对象拥有硬件资源，其他模块只能通过它的接口请求操作。比如串口对象维护发送队列、接收队列、错误计数和当前硬件状态。它不一定需要特权级，但已经有了内核思想：资源集中管理，状态对外封装，错误通过接口返回。

\begin{lstlisting}
struct UartDevice {
  RingBuffer rx;
  RingBuffer tx;
  uint32_t overrun_count;
  bool tx_busy;
};

bool uart_submit(UartDevice* dev, Span bytes);
bool uart_read(UartDevice* dev, byte* out);
\end{lstlisting}

从这里可以看到后续 fd、socket、设备文件的影子。用户不应该直接改设备寄存器，而应该持有一个句柄，请求拥有者执行操作。差别只是大型 OS 里这个拥有者运行在内核态，并且有硬件权限保护。

\subsection{从裸机模块边界到用户/内核边界}

裸机程序里也可以规定“只有 UART 模块能写 UART 寄存器”，但这只是团队约定。任何 C 文件只要拿到寄存器地址，仍然能直接写硬件。通用 OS 不能依赖所有程序守规矩，它必须让硬件禁止用户程序访问设备寄存器。

\begin{lstlisting}
bare-metal convention:
  app calls uart_submit()
  but app could still write UART_BASE directly

protected OS:
  user VA does not map UART MMIO
  user write to UART physical address triggers fault
  app must call write(fd, ...)
  kernel driver performs MMIO after permission checks
\end{lstlisting}

这就是模块化和隔离的区别。模块化让代码更清楚，但不能阻止恶意或错误代码；隔离由 CPU 特权级、页表权限和异常入口强制执行。OS 的很多复杂性，就是把裸机里的“请大家遵守约定”升级成“硬件保证不能绕过”。

\section{实践}

本章实践不需要复杂环境。拿任意一个单片机或模拟器，写三版程序：

\begin{enumerate}
\tightlist
\item 超级循环：固定顺序调用四个任务。
\item 周期任务表：每个任务有 period 和 next release。
\item 非阻塞状态机：把串口发送或传感器读取改成 step 函数。
\end{enumerate}

报告要记录每个任务最长运行时间、周期、deadline、是否可能阻塞、是否访问共享状态。对于每个任务，写出下面的表：

\begin{longtable}{p{0.2\linewidth}p{0.66\linewidth}}
\toprule
字段 & 内容 \\
\midrule
输入 & 读取哪些寄存器、队列或全局状态 \\
输出 & 写哪些寄存器、队列或全局状态 \\
时间预算 & 最坏执行时间和周期 \\
等待点 & 是否忙等硬件、锁或缓冲区 \\
失败 & 超时、溢出、校验失败、设备错误 \\
\bottomrule
\end{longtable}

\section{测试}

最低测试不是“灯会闪”，而是证明调度和状态机边界：

\begin{itemize}
\tightlist
\item 任一任务执行一次后必须返回主循环。
\item 周期任务的 \code{next\_release\_ms} 单调前进，不因时间溢出立刻失控。
\item 串口发送缓冲区满时，提交失败或施加背压，而不是覆盖旧数据。
\item 非阻塞状态机在硬件未 ready 时不改变发送进度。
\item 任一任务超时应被计数或记录，不能静默拖慢全系统。
\end{itemize}

这一章的验收结论是：操作系统不是凭空出现的。只要裸机程序里出现多个活动、共享资源、等待硬件、错误恢复和时间预算，就已经需要 OS 式的状态管理。

\section{源码级补充：真实 MCU 启动、UART 与链接脚本}

\subsection{x86-64 早期入口、\_start 与最小启动代码}

早期启动的第一份“内核源码”通常不是 C 函数，而是入口汇编和链接脚本。x86-64 上，UEFI/bootloader 或 EFI stub 把内核镜像放入内存，传入启动参数和内存地图；早期入口建立最小栈、清 BSS、准备早期页表和 IDT，然后才调用 \code{kernel\_main} 一类 C 入口。

\begin{lstlisting}
early_x86_64_entry:
  lea rsp, [rip + boot_stack_top]
  copy .data from flash LMA to sram VMA
  zero .bss
  build early page tables
  load provisional IDT
  call early_console_init
  call kernel_main
halt_loop:
  hlt
  jmp halt_loop
\end{lstlisting}

这里已经出现了 OS 启动的三个根概念：入口地址、内存布局和异常表。后面 IDT、内核入口和多核启动更复杂，但本质仍是“硬件和 bootloader 按固定规则跳到一段受信任代码，代码建立可运行环境”。

\subsection{VMA 与 LMA：为什么 data 要复制}

链接脚本里常见 VMA 和 LMA。VMA 是代码运行时使用的地址，LMA 是镜像中加载的位置。已初始化全局变量的初值保存在 Flash 中，但运行时变量要位于 SRAM，所以启动代码要从 LMA 复制到 VMA。

\begin{lstlisting}
MEMORY {
  FLASH (rx)  : ORIGIN = 0x08000000, LENGTH = 512K
  SRAM  (rwx) : ORIGIN = 0x20000000, LENGTH = 128K
}

SECTIONS {
  .text : { *(.text*) *(.rodata*) } > FLASH
  .data : AT(ADDR(.text) + SIZEOF(.text)) {
    __data_start = .;
    *(.data*)
    __data_end = .;
  } > SRAM
  __data_load = LOADADDR(.data);
  .bss : {
    __bss_start = .;
    *(.bss*) *(COMMON)
    __bss_end = .;
  } > SRAM
}
\end{lstlisting}

若 data 未复制，带初值的全局对象会从错误内容开始；若 BSS 未清零，锁、队列和指针的初始值不可预测。这个失败和大型 OS 的早期页表、per-CPU 区域未初始化一样，通常不会在第一条错误处崩溃，而是在后续随机位置表现出来。

\subsection{GPIO 和 UART 寄存器}

裸机驱动的核心是 MMIO。GPIO 通常有模式寄存器、输出数据寄存器、输入数据寄存器；UART 通常有状态寄存器、数据寄存器、波特率 divisor、控制寄存器。真实芯片位名不同，但结构类似。

\begin{lstlisting}
#define UART_SR   (UART_BASE + 0x00)
#define UART_DR   (UART_BASE + 0x04)
#define UART_BRR  (UART_BASE + 0x08)
#define UART_CR   (UART_BASE + 0x0c)
#define UART_TXE  (1u << 7)
#define UART_RXNE (1u << 5)

void uart_putc(char c) {
  uint32_t deadline = ticks + 10;
  while (!(read32(UART_SR) & UART_TXE)) {
    if (time_after(ticks, deadline)) {
      uart_errors.timeout++;
      return;
    }
  }
  write32(UART_DR, c);
}
\end{lstlisting}

这里的超时是工程边界。没有超时的驱动把硬件故障扩大成系统停顿。大型内核中的块设备超时、网络设备 reset、watchdog 都是同一思想：设备是异步状态机，不是可靠函数。

\subsection{轮询、中断和 CPU 占用率}

假设 UART 每秒收到 100 个字节，CPU 频率 100MHz。若主循环每 10 微秒轮询一次状态寄存器，一秒轮询 100000 次，其中只有 100 次有用；若每次轮询花 30 个周期，纯轮询检查就耗掉 300 万周期，约 3\% CPU。若轮询间隔变大，CPU 占用下降，但输入延迟和溢出风险上升。

\begin{longtable}{p{0.22\linewidth}p{0.34\linewidth}p{0.32\linewidth}}
\toprule
模式 & 优点 & 失败边界 \\
\midrule
忙轮询 & 实现最简单，延迟可控 & 浪费 CPU，慢操作会漏事件 \\
周期轮询 & 降低 CPU 消耗 & 延迟取决于周期，burst 易溢出 \\
中断驱动 & 空闲时不耗 CPU，事件到达即处理 & ISR 竞态、队列满、中断风暴 \\
DMA + 中断 & 大吞吐低 CPU & buffer 生命周期、cache、一致性复杂 \\
\bottomrule
\end{longtable}

这张表也适用于大 OS：epoll 减少无效轮询，NAPI 在高负载下把中断转为轮询，io\_uring 用共享队列减少系统调用往返。机制不同，目标都是在 CPU 占用、延迟和复杂度之间折中。

\subsection{MCU 上的 DMA：从串口接收到环形缓冲区}

即使在单片机上，DMA 也不是高级服务器才有的东西。典型场景是 UART 持续接收数据，CPU 不想每个字节都进中断。驱动可以准备一个环形缓冲区，让 DMA 控制器把 UART 接收数据直接写入 SRAM，半满或满时再中断 CPU。

\begin{lstlisting}
uart_rx_dma_setup():
  dma.src = UART_DATA_REGISTER
  dma.dst = rx_ring_buffer
  dma.len = RING_SIZE
  dma.mode = circular
  dma.enable_half_full_irq = true
  dma.enable_full_irq = true
  start_dma()

dma_irq():
  produced = compute_dma_write_position()
  publish_new_bytes_to_parser(produced)
\end{lstlisting}

这个例子把 DMA 的三条规则讲清楚。第一，buffer 在 DMA 运行期间不能随意改用途。第二，CPU 读取 DMA 写入的数据前，必须理解缓存一致性和写入位置。第三，completion 不一定表示全部 I/O 结束，也可能只是半缓冲区可消费。大型 OS 的网卡 RX ring 和块设备 completion，本质上是这个模型的扩展。

\subsection{真实代码阅读入口}

读真实 MCU 代码时，建议按下面顺序找：

\begin{enumerate}
\tightlist
\item 启动汇编或 startup 文件：向量表、reset handler、栈顶符号。
\item 链接脚本：Flash/SRAM 区域、\code{.text/.data/.bss}、堆栈布局。
\item HAL 或裸驱动：UART/GPIO/TIMER 寄存器定义和初始化序列。
\item 中断文件：IRQ handler 如何清中断、入队、唤醒主循环。
\item 示例 main：是否有阻塞等待、是否有超时、是否有错误计数。
\end{enumerate}

x86-64 教学内核常见文件可以包括 \filepath{boot/entry\_64.S}、\filepath{kernel/head64.c}、\filepath{linker.ld}、\filepath{drivers/early\_uart.c}。不要先读完整驱动框架，先把 UEFI/bootloader 交接到第一个字符输出的路径追通。
